% Partie : applied-modal-logics
% Chapitre : epistemic-logic
% Section : introduction

\documentclass[../../../include/open-logic-section]{subfiles}

\begin{document}

\olfileid[fr]{aml}{el}{int}

\olsection{Introduction}

De même que la logique modale étudie les \emph{propositions modales}
et les relations de conséquence logique entre elles, la logique
épistémique étudie les \emph{propositions épistémiques} et les relations
de conséquence logique entre elles. Les connecteurs unaires reçoivent
ici une interprétation épistémique ou doxastique, destinée à représenter
la connaissance et la croyance, au lieu de représenter la possibilité
et la nécessité. Nous pourrions, par exemple, vouloir exprimer les
énoncés suivants :

\begin{enumerate}
\item Richard sait que Calgary est en Alberta.
\item Audrey pense qu'il est possible qu'un chien soit sur le canapé.
\item Richard sait qu'Audrey sait que son cours a lieu le mardi.
\item Tout le monde sait qu'une année compte 12 mois.
\end{enumerate}

On fait souvent remonter la logique épistémique contemporaine à
\emph{Knowledge and Belief}, publié par Jaakko Hintikka en 1962\footnote{L’orthographe Jaako de la source est corrigée en Jaakko, conformément au \href{https://data.nationallibrary.fi/bib/me/W00031358300}{catalogue de la Bibliothèque nationale de Finlande}.}, à une
époque où la sémantique des mondes possibles se répandait en logique.
Les logiques épistémiques utilisent, en effet, pour l'essentiel les
mêmes outils sémantiques que les autres logiques modales, mais elles
les interprètent différemment. Le principal changement porte sur ce
que représente la \emph{relation d'accessibilité} : une forme de
\emph{possibilité épistémique}. Nous verrons que la notion épistémique
que nous cherchons à représenter influe sur les contraintes à imposer
à cette relation. Nous verrons aussi ce que devient la théorie de la
correspondance lorsqu'on lui donne une interprétation épistémique.
Les exemples précédents font intervenir deux agents, Richard et
Audrey, ainsi que les liens entre leurs connaissances respectives.
Les logiques épistémiques que nous étudierons sont des logiques à
plusieurs agents, qui permettent d'exprimer ces liens. Une logique
épistémique à un seul agent ne porte, quant à elle, que sur ce qu'un
individu sait ou croit.

\end{document}
